\section{All auxiliary arities of a guarded native block}
\label{sec:136-native-all-arities}

\paragraph{Archive and exact scope.}
This concerns cote136, physical p.~117/footer116/scan117,
physical p.~119/footer118/scan119, and the common-source question
physical p.~121/footer120/scan121. It strengthens the guarded
native residue statement, using the checked linear factorization
in ARCH-136-HIGH-NATIVE-RESIDUE. It concerns a homogeneous native
block and its specified RH1 iteration. It is not an arbitrary
old-algebra pushout or a theorem about later horizontal repairs.

Let \(p\) be prime, \(R=\mathbb Z_{(p)}\),
\(S=\Gamma_R(T)\), and begin with a free odd unit disk
\(a,b\), where \(|a|=q\geq1\) is odd,
\(da=b\), and both generators have weight \(m>0\).
Give them auxiliary arity one. Put
\[
 e=a\gamma_{p-1}(b),\quad z=\gamma_p(b),\quad
 D=p(q+1),\quad de=pz,\quad v=pm.
                                                               \tag{136.NA1}
\]
Adjoin the full coherent tail presentation on a free resolution
of \(I_M\), with all its \(U_i,V_i\) generators assigned
auxiliary arity \(p\), including both bottom relations.
Call the resulting native block \(A\). Its arity-\(p\)
linear complex \(K\), including \(e,z\), is the truncated
resolution of \(S/(p,I_M)\langle v\rangle\) in forms
zero through \(D\). This follows by the same exact sequence
as (136.NM4), now with \(p\) instead of two. The coefficient
tail is integer-saturated, so the connecting inclusion
\(I_M/pI_M\to S/pS\) remains injective.

Use the homogeneous simplicial model on
\(\Delta[n]/\operatorname{sk}_1\Delta[n]\), or its
representable family version, so the native positive-arity
part is zero on the first skeleton. Suppose maps to a finite
family of admissible targets are given with that relative
property, and assume
\[
                    pm\geq\max_j W_{C_j}(D)
                                                               \tag{136.NA2}
\]
with the native-weight bound of ARCH-136-HIGH-NATIVE-RESIDUE.

\paragraph{ARCH-136-NATIVE-ALL-ARITIES.}
Every degree-three Moore pair in every positive auxiliary arity
of this native model has zero individual graph residue in all
the targets. The same holds after the homogeneous RH1 iteration,
including its filtered colimit. The result does not require
annihilation of source cohomology by a coefficient tail.

\begin{proof}
First identify the entire old arity-zero residue algebra,
where ``zero residue'' means auxiliary arity divisible by \(p\).
It is exactly
\[
              B=\Gamma_S(z)\otimes\Lambda_S(e).
                                                               \tag{136.NA3}
\]
Indeed the coefficient of \(\gamma_{pr}(b)\) in
\(\gamma_r(z)\), and the coefficient of
\(a\gamma_{pr-1}(b)\) in \(e\gamma_{r-1}(z)\),
are both
\[
                      c_r=\frac{(pr)!}{r!(p!)^r}.
\]
Legendre's factorial valuation gives \(v_p(c_r)=0\), so
these are units in \(R\). The displayed generators therefore
give every basis element, with its correct differential.

For each other auxiliary residue \(j\ne0\pmod p\),
the old native module has the contraction
\[
 h(\gamma_s(b))=\frac{a\gamma_{s-1}(b)}{s},\qquad
 h(a\gamma_{s-1}(b))=0,
                          \quad s\equiv j\pmod p.
                                                               \tag{136.NA4}
\]
Here \(s\) is a unit and \(dh+hd=1\).
It is a graded module homotopy over \(B\).
For the generator \(z\), the necessary identity is
\[
 \frac{\binom{s+p}{p}}{s+p}
                 =\frac{\binom{s+p-1}{p}}{s}.
\]
The same identity with \(p\) replaced by \(pr\) proves
compatibility with every divided power of \(z\), including
its scalar \(c_r\). For \(e\), both relevant compositions
are zero because of the strict square-zero rule for \(a\).
Thus (136.NA4) is genuinely \(B\)-linear, with the usual
graded sign, not merely coefficient-linear after rationalization.

All coherent parameters have arity divisible by \(p\), and
their differential lies in those parameters together with
\(e,z\). The arity-divisible part of \(A\) is consequently
the full algebra \(\Gamma_{\mathrm{gr}}(K)\), with no
restriction to quadratic terms. Every other residue component
is the corresponding old \(B\)-module tensored with that
algebra over \(B\). Extend (136.NA4) by module linearity.
The terms differentiating the new algebra cancel in
\(dh+hd\), so this is still a contraction. It preserves
weight and each auxiliary arity. No flatness is needed.

In a simplex with several independent native basis disks,
use one arity coordinate for each such disk. The component
whose coordinates are all divisible by \(p\) is
\(\Gamma_{\mathrm{gr}}(\bigoplus K_i)\). Every other
component contracts by the preceding homotopy on its first
nonzero residue coordinate, with the usual tensor sign.
It has no zeroforms. These are pointwise decompositions;
faces may identify native basis disks. The all-zero-residue
subalgebra is nevertheless simplicial, since such identifications
add multiples of \(p\). All subsequent uses of a divisible
component mean this componentwise condition when there is
more than one native disk. The uniform linear homotopies and
free disk factorization apply to the direct sum, one free
parameter component at a time.

For the divisible component, apply
ARCH-136-HIGH-NATIVE-RESIDUE to the linear map
\(K\to C_j(\Delta[n],\operatorname{sk}_1\Delta[n])\).
Condition (136.NA2) supplies linear homotopies in this kernel.
The universal map \(\Gamma_{\mathrm{gr}}(K)\to E\)
constructed in ARCH-136-RELATIVE-DISK-FACTORIZATION is the
same for every target; only the evaluation of its free disk
generators differs. The free disk algebra \(E\) has
\(H^0=S\), \(H^1=H^2=H^3=0\) in every weight and
positive auxiliary arity. Its positive-arity part maps into
the graph kernel in every target. Therefore all the source
Moore pairs in this component have primitives in \(E\),
whose target images have graph values zero. For a nondivisible
component, its actual contraction already makes its Moore
cone acyclic, so the same conclusion follows from source
primitives. This proves the assertion before RH1.

We verify the RH1 continuation, rather than assuming that
its new classes came from the old algebra. At every stage,
the nondivisible components remain contracted: all zeroforms
have divisible auxiliary arity, and the contracted components
have no one-cycles, since they have no zeroforms either.
RH1 thus adjoins only zeroforms of divisible auxiliary arity.
The module contractions extend by base change over the
enlarged divisible algebra. All identifying PD relations are
generated by zeroforms in that algebra; the scalar divided-power
rule keeps their positive-form PD consequences in the same
generated module ideal. Thus this is actual module base change,
with the same differential
cancellation. This holds at successor stages and at the
filtered colimit.

It remains to extend the factorization of the divisible
component through these zeroform repairs. This can be checked
at each fixed simplex; a simplicial structure on the auxiliary
factor \(E\) is not needed. In any positive auxiliary arity,
the RH1 replacement module at that simplex is precisely the
module of old one-cycles in the weights being repaired.
Indeed normalization commutes with kernels, the normalized
degree-zero one-form module is zero because the source has
simplex zero exactly \(S\), and Dold--Kan reconstructs the
direct sum of the normalized cycle modules. In lower weights
RH1 is the identity.

Map each such cycle to its image in \(E^1\) and take its
unique positive-arity zeroform primitive there. Existence uses
\(H^1(E)=0\); uniqueness uses the absence of positive-arity
\(H^0(E)\). This assignment is linear and agrees with the
old zeroforms on all identifying relations. The free PD
pushout therefore extends the map to \(E\). It continues
to agree with each retained target map: both images of a new
zeroform are primitives of the same target one-cycle; both
vanish at vertices, while the target's closed zeroforms in
the fixed repair range are exactly its coefficients. Their
difference is consequently zero. The evaluation from \(E\)
still maps positive arity into the graph kernel. The same
unique primitive construction extends the factorization at
every stage, compatibly, and hence at the colimit.

Every new positive-arity degree-two cone pair in the divisible
part therefore still factors through this same \(E\), whose
degree-two cone cohomology is zero. The nondivisible parts
remain contracted as proved above. All its target primitive
pairs have zero graph restrictions. This proves the complete
stated RH1 assertion without assuming any annihilator for
the nonlinear source cohomology.
\end{proof}

\paragraph{Limits of this conclusion.}
The native-weight condition (136.NA2), the homogeneous block
presentation, and its graph-zero target maps are genuine
hypotheses. A map from a general old algebra need not factor
through these blocks; arbitrary subsequent horizontal or
positive-form attachments are not covered. For an integral
all-prime construction, the bounds
\(pm\geq W_C(p(q+1))\) are not known to follow from one
finite lower bound on \(m\). This result therefore does
not assert integral connectedness or a complete compatible
source construction.

\paragraph{Verification status.}
The author has read the complete saved proof, including the
unit factorial coefficients, full module contraction, free
factorization, pointwise Dold--Kan interpretation of RH1,
unique primitives, target compatibility at vertices, and
the exact native-weight and continuation restrictions.
The parent supplied the full residue-module contraction idea;
the proof uses the independently checked factorization of
reviewer1348. Root completed a whole-proof PASS with contributor scope.
Reviewer1348 has now completed a whole saved-proof read and
communicated PASS on 2026-11-22, including the pointwise RH1
factorization and the multiple-block residue coordinates.
